\section{Trabajo relacionado}
\label{sec:related_work}

\noindent\textbf{LLMs para Video Understanding.} 
En los últimos años, ha habido un fuerte aumento del interés por aprovechar Large Language Models (LLMs) para video understanding, lo que ha llevado al desarrollo de varios enfoques innovadores~\citep{video_chatgpt, llava_next_video, llava_onevision}. Estos modelos suelen usar un Vision Encoder para extraer características de video, seguido de una etapa de mapeo con Linear Projection, MLP o Q-Former~\citep{blip2}. Las características mapeadas se combinan con datos textuales y se introducen en large language models (LLMs) para generar una salida de texto. Estos modelos tienen arquitecturas relativamente simples, requieren menos datos de entrenamiento y recursos computacionales, y aun así logran un rendimiento sólido en benchmarks de short video understanding~\citep{msrvtt_qa, next_qa, li2024mvbench}. Sin embargo, emplean sparse sampling o técnicas de token compression para reducir el número de tokens, lo que puede producir una pérdida significativa de información al tratar con videos más largos o con mayor riqueza de contenido. Como resultado, no son adecuados para long video understanding ni para streaming video understanding.

\noindent\textbf{Long Video Understanding.} 
Un desafío central en long video understanding es comprimir eficazmente la información de videos extensos. Muchos enfoques usan el lenguaje como puente, condensando videos en dense captions~\citep{llovi,video_recap,streaming_dvc}. Aunque esto logra buenos resultados en algunos casos, comprimir contenido de video en texto suele conducir a la pérdida de detalles visuales cruciales. Además, como enfoque pionero para streaming video understanding, VideoLLM-Online~\citep{videollm_online} emplea una metodología centrada en los datos al intercalar video y texto durante el entrenamiento. En contraste, nuestro enfoque es training-free, lo que permite una integración transparente con varios Video-LLMs existentes para extender sus capacidades de StreamingVQA. Adicionalmente, VideoLLM-Online conserva solo un único token por frame para manejar videos largos, lo que puede provocar pérdida de información visual. Nuestro método preserva la información visual completa y aprovecha In-Context KV-Cache Retrieval para mejorar la eficiencia.

Otra línea de investigación se centra en comprimir videos largos en un memory bank~\citep{wu2019long, wu2022memvit, wang2023memory}.
MC-ViT~\citep{mc_vit} adapta video transformers preentrenados existentes mediante fine-tuning para que atiendan a memorias visuales condensadas. Se relaciona estrechamente con los métodos de token-pruning, merging y memory-based video understanding. En comparación, proponemos un método training-free específicamente orientado a la tarea de StreamingVQA. Incorporar MC-ViT en la tarea de StreamingVQA podría ser una vía interesante para investigaciones futuras, y reconocemos su potencial en este dominio.
Este enfoque se ha integrado en Video-LLMs para streaming video understanding, como se muestra en trabajos como VideoStreaming~\citep{videostreaming} y Flash-VStream~\citep{flashvstream}. Estos métodos actualizan dinámicamente la memoria durante el procesamiento de video y la utilizan para tareas downstream.
A pesar de su innovación, una limitación importante de estos métodos es que no consideran la longitud del video ni la densidad de información, especialmente cuando usan un tamaño de memoria fijo. Por ejemplo, Flash-VStream comprime tanto clips de 10 segundos como videos de una hora en los mismos 681 tokens.
Además, estos métodos carecen de interpretabilidad, lo que dificulta determinar cuánta información se está comprimiendo en la memoria o si la información de video relevante se está recuperando correctamente durante las tareas downstream.

En la búsqueda de una mayor interpretabilidad en long video understanding, métodos como GroundVQA~\citep{di2023groundvqa} y GeLM~\citep{chen2024groundedmultihopvideoqalongform} proponen localizar clips de video relevantes mientras responden a consultas del usuario. Inspirado en ellos, este trabajo evita condensar en exceso la información de video. Al aprovechar las capacidades causales de los Video-LLMs, preserva la Video KV-Cache completa, permitiendo recuperar la información relevante cuando se la requiere. Esta estrategia mitiga eficazmente la pérdida sustancial de contenido de video y, al mismo tiempo, mejora la interpretabilidad.

\noindent\textbf{Manejo de long context para LLMs.} 
El manejo de secuencias de texto largas en LLMs ha sido un desafío importante debido a los altos costos computacionales y de memoria, lo que ha impuesto restricciones de entrenamiento sobre secuencias más cortas.
Técnicas como StreamingLLM~\citep{streamingllm} y LM-Infinite~\citep{lm_infinite} utilizan atención de ventana deslizante para procesar secuencias largas de forma incremental, pero descartan tokens lejanos, limitando la capacidad del modelo para capturar dependencias de largo alcance. Enfoques recientes~\citep{infllm,quickllama,em_llm} abordan esto almacenando y recuperando KV-Caches previamente calculadas, habilitando un mejor recall de contextos distantes.

\noindent\textbf{Retrieval-Augmented Generation.} 
Retrieval-augmented generation (RAG) combina mecanismos de retrieval con modelos generativos para mejorar el rendimiento en diversas tareas de NLP al incorporar conocimiento externo~\citep{guu2020retrieval,lewis2020retrieval,borgeaud2022improving} y mejorar el rendimiento en tareas vision-language~\citep{xu2024retrieval}. In-context retrieval, propuesto recientemente para manejar entradas largas~\citep{ram2023context}, recupera información del propio documento de entrada en lugar de una base de conocimiento externa. In-context KV-Cache retrieval mejora aún más la eficiencia al pre-encodificar documentos largos, evitar codificaciones redundantes y aprovechar las capacidades de retrieval del propio LLM para un rendimiento más rápido y efectivo.
